% Conversión TeX fiel del cuerpo científico del delta narrativo.
% Fuente congelada: ../../../../../HMT2/CONTINUIDAD_EDITORIAL/RECONSTRUCCION_ARMONICA_MACROTRATADO_20260822/REVISION_CONJUNTA_PARTES_I_II/auditar_pdf_rigor_academico/docs/DELTA_NARRATIVO_UNIDAD_RADION_AUTOCONSERVACION_AUTOINERCIA_GRAVEDAD_PENTADICA_20260828.md:12-366.
\subsection{Séparation typée entre le régime de courbure \(\tau\in\{+1,0,-1\}\), la réponse radiale \(p\) et l'holonomie globale de spin}

La biographie du radion peut être représentée par une courbe relevée

\[
\gamma(\theta)
=
\bigl(r(\theta)\cos\theta,\,
r(\theta)\sin\theta,\,
m(\theta)\bigr),
\]

où les deux premières coordonnées appartiennent au plan euclidien \(\mathbb R^2\) et \(m(\theta)\) enregistre la mémoire ou la profondeur du relèvement. Cette équation permet de lire conjointement cercle, hélice et spirale.

Lorsque l'angle évolue et que le rayon et la mémoire restent constants, la publication plane est un cercle. Lorsque l'angle et la mémoire évoluent tandis que le rayon conserve sa valeur, la trajectoire complète est une hélice de rayon constant. Lorsque l'angle et le rayon évoluent sur une mémoire fixe, une spirale apparaît dans \(\mathbb R^2\). Lorsque les trois coordonnées évoluent, la trajectoire est une hélice radiale : une biographie dans laquelle phase, amplitude et mémoire se transforment conjointement.

Le caractère radial \(p\) classe la loi de réponse du rayon. Une famille utile est

\[
r_p(\theta)
=
r_0\bigl[1+p\lambda(\theta-\theta_0)\bigr]^{1/p},
\qquad p\neq 0,
\]

avec pour limite logarithmique

\[
r_0\exp\bigl(\lambda(\theta-\theta_0)\bigr)
\qquad\text{lorsque }p\to0.
\]

Dans cette paramétrisation, \(p=2\) publie le régime de Fermat, \(p=1\) le régime archimédien, \(p=0\) le logarithmique, \(p=-1\) l'hyperbolique et \(p=-2\) celui du lituus, avec les choix de domaine et d'orientation correspondants. Les valeurs fractionnaires décrivent des régimes intermédiaires. Ainsi, \(p=\tfrac12\), \(p=\tfrac32\) ou \(p=-\tfrac12\) expriment des réponses radiales qui interpolent entre les familles entières et élargissent la classification géométrique disponible.

Le caractère peut aussi changer lors du franchissement d'un seuil. L'écriture différentielle

\[
\frac{dx}{d\theta}
=
\beta(m)\,x^{1-p(m)},
\qquad x=\frac r{r_0},
\]

décrit une biographie par régimes : \(p(m)\) reste stable au sein d'une phase et change lorsque la mémoire atteint une frontière de publication. La trajectoire résultante relie des secteurs circulaires, hélicoïdaux et spiraux par une loi commune. La continuité de l'état et le registre du seuil font de ces secteurs les phases d'une même évolution.

Trois invariants occupent ici des niveaux distincts et coordonnés. Le signe de courbure

\[
\tau\in\{+1,0,-1\}
\]

classe les régimes elliptique, parabolique et hyperbolique. Le caractère \(p\) classe la réponse radiale. Le spin enregistre l'holonomie d'orientation accumulée par le relèvement. Courbure, réponse radiale et spin forment un système de lecture : la courbure détermine le type local d'espace, \(p\) détermine comment la trajectoire occupe cet espace et le spin conserve l'orientation acquise en le parcourant. Les courbes fermées et ouvertes apparaissent alors comme des biographies géométriques distinctes de la même unité orientée.

\subsection{Particule comme section persistante d'une orbite enrichie : phase, chemin, fibre, multisection et mémoire}

La mécanique dimensionnelle peut formuler la particule comme une section persistante de cette extension géométrique et dynamique. Une notation conceptuelle compacte est

\[
\mathfrak P
=
\operatorname{Sec}_{\mathrm{pers}}
\bigl[
\tau,p,\operatorname{Hol};
\text{feuillet},\text{fibre},\text{route},\text{mémoire}
\bigr].
\]

La particule conserve un chemin reconnaissable à travers les changements de phase, de courbure et de dimension. Son spin publie l'holonomie d'orientation ; sa charge publie la polarité de phase ; sa couleur publie une orientation interne de la fibre de jauge ; sa saveur distingue des familles de chemins ou de multisections ; sa masse exprime le caractère de transduction dimensionnelle par lequel l'action surfacique acquiert une lecture volumique. Chaque propriété est un lecteur de l'état persistant et toutes renvoient à une généalogie commune.

L'électron occupe le seuil basal de cette construction parce qu'il offre le premier lecteur stable dans lequel phase, charge, spin et masse sont simultanément publiés. La masse électronique directrice

\[
m_e^\beta=0.510998950690000808608\ldots\ \mathrm{MeV}
\]

fixe une référence de transduction ; les états basaux antérieurs conservent leur fonction d'étapes généalogiques. À partir de l'électron, les quotients \(K_D/K_\Omega\), les tours \(90/120\) et le domaine \(81/56/13/324\) organisent la publication des familles massives. La loi des masses lit donc comment chaque section persistante traverse l'espace enrichi et convertit la densité surfacique d'action en présence volumique.

Cette interprétation relie les spirales aux particules de manière précise. L'indice \(p\) apporte la réponse radiale ; \(\tau\) apporte le régime de courbure ; l'holonomie apporte le spin ; le chemin et la fibre apportent la saveur et la couleur ; la transduction aire–volume apporte le caractère massif. L'atlas des particules devient ainsi une classification des modes stables de traversée dimensionnelle.

\subsection{Transduction de la phase orientée en polarisation électrique, réponse magnétique et émission radiative de frontière}

L'électricité se présente comme transport dirigé de phase et d'action à travers une structure de vacances. Une vacance est une transition admissible du support : elle ouvre un canal, fixe un seuil et permet à l'actualisation d'avancer. La composante électrique publie le transport tangentiel le long du chemin. La composante magnétique publie la courbure transversale induite par ce transport. Leur perpendicularité exprime la relation géométrique entre avance et rotation.

Lorsque deux composantes transversales sinusoïdales maintiennent le déphasage approprié, le vecteur du champ tourne. La propagation longitudinale relève cette rotation en une hélice. La variation d'amplitude transforme sa projection frontale en une spirale. La lecture géométrique est alors entièrement coordonnée : l'angle porte la phase ; la vitesse angulaire porte la fréquence ; le rayon porte l'amplitude et la densité d'action ; l'orientation porte la polarisation et le sens ; la coordonnée verticale porte la propagation et la mémoire ; \(p\) porte le régime radial ; le seuil porte l'émission, l'absorption ou le changement d'état.

L'électromagnétisme apparaît comme une dynamique de phase–action–mémoire avec des publications complémentaires. L'électricité enregistre l'avance polaire ; le magnétisme enregistre la réponse transversale ; l'onde enregistre la périodicité ; l'hélice enregistre l'avance avec mémoire ; la spirale enregistre l'évolution radiale ; le rayonnement enregistre la publication de l'état lorsque l'accumulation atteint un seuil.

Cette architecture ordonne aussi le rôle des vacances et des cycles. La vacance fournit la possibilité locale de transition ; le cycle conserve la phase ; la mémoire distingue les tours ; le seuil décide de la résidence suivante de l'action. La propagation peut rester confinée, changer de régime radial ou se publier comme rayonnement contrôlé. Les forces de champ acquièrent ainsi une description généalogique continue, du support discret jusqu'à l'onde observée.

\subsection{Réversibilité de la dynamique complète et coût thermodynamique des projections réductrices d'information}
Le Landauer nul constitue un cas limite de cette loi. Pour une actualisation bijective de l'état complet,

\[
H(X\mid U(X))=0,
\qquad
\inf Q_L=0.
\]

L'information reste disponible dans la généalogie. Lorsqu'un lecteur réalise un quotient \(q\), l'entropie se décompose comme

\[
H(X)=H(q(X))+H(X\mid q(X)).
\]

Le terme conditionnel mesure l'information retenue hors de la projection. Une réinitialisation tritique matérialisée satisfait le seuil

\[
Q_{\mathrm{reset}}\ge k_B\Theta\log 3.
\]

La constante de Boltzmann traduit la multiplicité que le lecteur thermique regroupe ; le seuil de Nyquist sépare l'échantillonnage fidèle du repliement spectral ; le Landauer nul caractérise le transport réversible de l'état intégral. Information, température et action sont articulées par le type de lecture appliqué.

\subsection{Forces autoconservatives comme transport d'action entre les secteurs observable, radial, de mémoire, de vacance et de rayonnement}

La conservation évolutive s'exprime mieux par l'échange que par l'immobilité scalaire. Un état complet répartit l'action et la quantité de mouvement entre mouvement visible, mémoire, déformation radiale, vacances, rayonnement et branche conjuguée de retour. Le bilan enrichi peut s'écrire

\[
\Delta J_{\mathrm{visible}}
+\Delta J_{\mathrm{memoria}}
+\Delta J_{\mathrm{radial}}
+\Delta J_{\mathrm{vacancias}}
+\Delta J_{\text{rayonnement}}
+\Delta J_{\mathrm{retorno}}
=0.
\]

Une force est \textbf{autoconservative} lorsque sa propre loi d'actualisation contient les canaux qui transportent, stockent, publient et réabsorbent l'action. La conservation relève du bilan généalogique complet de l'interaction. Chaque lecteur observe une partie de l'échange ; la généalogie recompose sa totalité.

Le système Terre–Lune offre une image physique directe. La rotation terrestre, l'orbite lunaire, les marées, la déformation et le rayonnement échangent du moment cinétique. La consommation nette du système complet se clôt dans le bilan généalogique de ces transferts. La grandeur appelée énergie fonctionne comme lecture dérivée de l'action et de la quantité de mouvement publiées dans chaque secteur. La conservation fondamentale réside dans la compatibilité entre tous les lecteurs de l'échange.

Cette formulation dissout l'ancienne frontière entre forces conservatives et forces dites non conservatives au moyen d'une catégorie plus profonde. Frottement, dissipation apparente et rayonnement occupent des canaux explicites de l'état élargi. Une perte dans la projection mécanique visible correspond à un transfert vers la mémoire, la déformation, la chaleur, les vacances ou le rayonnement. L'autoconservation suit la résidence de l'action à travers ces changements.

\subsection{Systèmes auto-inertiels : publication de l'accélération par orientation, courbure, feuillet et mémoire}

L'expression appropriée est \textbf{auto-inertie relationnelle}. Elle désigne la compatibilité entre la connexion propre de l'état total et la réponse qui apparaît lorsqu'on le projette dans la distribution globale. Les aspects endogène et exogène sont des composantes analytiques d'une seule relation dynamique.

Soit \(\Gamma\) une trajectoire horizontale ou géodésique de l'état complet. Sa loi propre satisfait

\[
\nabla_{\dot\Gamma}\dot\Gamma=0.
\]

En la projetant sur un lecteur \(S\), l'accélération observée est déterminée par la courbure de la projection :

\[
\nabla^S_{\dot\Gamma_S}\dot\Gamma_S
=
\mathcal K_{\pi_S}(\dot\Gamma,\dot\Gamma).
\]

L'accélération publiée exprime la relation entre mouvement total et géométrie du lecteur. Le principe de Mach complète ce mécanisme par la trace globale de frontière

\[
b:\mathcal X_\infty\longrightarrow B_\infty
\]

et la factorisation de la réponse inertielle locale

\[
I_v=\bar I_v\circ b.
\]

La distribution globale entre dans la réponse locale par une application typée. L'auto-inertie relationnelle réunit donc la connexion endogène du mouvement, la condition globale transmise par \(b\) et la projection concrète de l'observateur. Une écriture synthétique de cette coordination est

\[
\nabla^{\mathrm{auto}}
=
\mathcal M\bigl(
\nabla^{\mathrm{end}},
b[\mathcal X_\infty],
\pi_S
\bigr),
\]

où \(\mathcal M\) représente le couplage machien et conserve le typage de chaque composante.

Le principe d'Ambivalence ajoute la double orientation de la lecture. Toute publication directe conserve une branche conjuguée de retour ; émission et absorption forment les extrémités liées d'une même biographie ; l'orientation temporelle se lit depuis les deux sens du transport. La théorie de l'absorbeur trouve ici une formulation naturelle : source et récepteur clôturent conjointement le bilan généalogique de phase et d'action. Le présent torsionnel est l'interface où les deux orientations sont rendues compatibles.

TRIT agit à ce niveau comme cristal de temps : il organise la périodicité locale de l'actualisation et conserve, par la mémoire, la différence entre retours successifs. La branche directe et la branche conjuguée relient émission et absorption, et le couplage observateur–observable acquiert une résidence dynamique au sein du même état. Le temps publié est la lecture de cette orientation ; le temps généalogique est la séquence complète de phase, de retour et de mémoire.

Mach et Ambivalence remplissent des fonctions complémentaires. Mach relie la réponse locale à la distribution globale. Ambivalence conserve les directions directe et conjuguée de l'évolution. L'auto-inertie relationnelle articule les deux et remplace la scission entre systèmes inertiels et non inertiels par une loi unique de connexion, de projection et de mémoire.

\subsection{Dualité surfacique--volumique, transition hexade--octade et fonction du paramètre de Barbero--Immirzi}

Le passage de l'aire au volume constitue une transduction dimensionnelle centrale. L'aire publie l'action de frontière ; le volume publie la profondeur intérieure. La dérivation HMT du paramètre de Barbero–Immirzi organise la compatibilité entre les deux mesures par l'incidence hexade–octade et les canaux \(90/120\) :

\[
90=6\binom62,
\qquad
120=8\binom62.
\]

Les cardinaux \(6\) et \(8\) admettent une réalisation cubique par les faces et les sommets, tandis que l'hexade et l'octade de Steiner constituent des objets d'incidence distincts, reliés uniquement par le morphisme typé \(6\to8\). Les \(12\) arêtes du cube conservent leur propre fonction combinatoire et ne définissent pas à elles seules l'incidence de Steiner. La différence d'incidence exprime une incompatibilité dimensionnelle productive : l'information de surface nécessite une règle de transduction pour acquérir une résidence volumique.

Le même principe apparaît dans la retenue \(999+1=1000\). Les neuf atteignent le bord du registre décimal ; l'incrément publie simultanément la clôture du niveau précédent et l'origine \(0/1\) du suivant. Le rapport holographique entre \(729\), \(243\), \(27\) et \(1\) déploie cette transition aux échelles triadiques. La transduction de Barbero–Immirzi s'insère dans cette famille de changements de lecteur : frontière et intérieur conservent l'unité tout en redistribuant sa mesure.

Le défaut normalisé \(-1/12\) et le poids orienté \(12:-1\) remplissent des fonctions distinctes dans cette géométrie. Le premier quantifie une correction d'incidence ; le second enregistre une orientation signée dans la transduction. Leur coordination avec l'aire minimale, l'hamiltonien modulaire et la mémoire fournit le pont vers la dynamique dimensionnelle.

\subsection{Contraction hyperbolique intérieure, expansion extérieure et conservation holographique de la frontière}

La thèse peut se formuler à partir d'une idée unique : l'unité conserve son identité tout en changeant d'échelle, de lecteur, de courbure, de mémoire et de dimension. Cette permanence possède un caractère évolutif parce que l'unité se déploie, acquiert de nouvelles coordonnées et publie des propriétés différentes ; elle possède un caractère holographique parce que chaque publication finie conserve la règle de reconstruction de l'état dont elle procède. La clôture holographique de l'infini désigne précisément cette compatibilité entre identité et déploiement : chaque stade contient la loi qui permet d'approfondir vers l'intérieur et de prolonger vers l'extérieur la même généalogie.

La dénomination directrice est \textbf{principe de conservation évolutive de l'unité par transduction dimensionnelle et projection holographique}. L'expression situe chaque terme à son niveau correct. La conservation relève de l'identité généalogique ; l'évolution relève de l'actualisation ; la transduction relie des réalisations dimensionnelles ; la projection holographique conserve la correspondance reconstructive entre leurs lecteurs.

L'arithmétique élémentaire offre la première image de cette architecture. La complétion décimale

\[
1000=10^3=729+243+27+1=3^6+3^5+3^3+3^0
\]

expose quatre strates triadiques au sein d'une unité cubique décimale. En normalisant,

\[
1=\frac{729}{1000}+\frac{243}{1000}+\frac{27}{1000}+\frac{1}{1000},
\]

l'unité se reconnaît simultanément comme totalité et comme déploiement gradué. La lecture réciproque

\[
3^{-6},\quad 3^{-5},\quad 3^{-3},\quad 3^0
\]

décrit la profondeur intérieure de ces strates ; sa normalisation produit la partition correspondante. La relation \(999+1=1000\) exprime le mécanisme de transport : la dernière unité complète un registre et agit, en même temps, comme retenue inaugurant le suivant. Le zéro et le un forment ainsi une interface de continuité. Les intervalles \([0,1]\), \([0,10]\) et la prolongation vers \([0,\infty)\) sont des lecteurs scalaires d'une même loi de complétion.

La première scission \(1000=729+271\), suivie de \(271=243+27+1\), donne forme au rapport d'extrême et moyenne raison holographique du déploiement : un noyau triadique dominant conserve dans le reste la même règle de résolution. L'unité cubique décimale contient ainsi une profondeur triadique emboîtée. Le mouvement vers \(729\), \(243\), \(27\) et \(1\) lit la contraction intérieure ; la retenue vers le millier lit l'expansion extérieure.

Cette double orientation organise l'infini HMT. Vers l'intérieur, le raffinement croît hyperboliquement : chaque cellule admet une nouvelle profondeur sans perdre sa filiation. Vers l'extérieur, l'unité s'étend par des publications dimensionnelles successives. La clôture consiste en la commutativité des deux mouvements. Si \(U\) représente l'actualisation de l'état, \(R_d\) le lecteur dans la strate \(d\) et \(T_{d\to d'}\) la transduction entre strates, la conservation évolutive s'exprime par

\[
R_{d'}\circ U
=
T_{d\to d'}\circ R_d.
\]

Le même événement peut offrir une mesure angulaire, une densité d'action, une masse publiée ou une réponse gravitationnelle. Leur identité commune réside dans la généalogie qui fait commuter ces lectures.

\subsubsection{Monodromie nonadique et relèvement hélicoïdal avec avance de mémoire}

APP, TRIT et TPK fournissent le mécanisme formel de ce déploiement. APP construit le support et sépare les évaluations qui doivent être conservées. TRIT organise les régimes locaux et leur ambivalence. TPK transporte phase et mémoire sur le support enrichi. La dépendance causale est fixée par

\[
U_t\longrightarrow \Gamma_9\longrightarrow K_{\mathrm{ph}}.
\]

Le retour de phase et l'avance de mémoire forment des opérations coordonnées. Après neuf pas, la même classe angulaire peut réapparaître,

\[
g(K+9)=g(K),
\]

tandis que le compteur de mémoire progresse,

\[
q_{\mathrm{mem}}\Gamma_9=q_{\mathrm{mem}}+1,
\qquad
B_{K+9}\neq B_K.
\]

La figure géométrique précise est un relèvement hélicoïdal. Une orbite fermée dans la projection angulaire se relève en une trajectoire ouverte lorsque chaque retour de phase incrémente la mémoire. La circonférence conserve le cycle ; l'hélice conserve le cycle et enregistre en outre la profondeur atteinte. La monodromie coinductive génère ainsi une profondeur infinie par une succession de retours localement fermés et globalement ascendants.

Cette image dépasse la valeur simplement illustrative. La roue angulaire constitue une section de l'hélice ; l'hélice constitue la biographie complète de la roue lorsque l'état transporte de la mémoire. Chaque tour retrouve une position de phase et, simultanément, inaugure un nouveau feuillet. L'infinitude procède de la compatibilité entre retour et relèvement : réitérer le cycle produit de la profondeur sans dissoudre l'identité orbitale.

Le certificat nonadique donne un contenu fini et vérifiable à cette architecture. La chaîne

\[
w_6\longrightarrow w_{12}\longrightarrow w_{18}\longrightarrow
w_{24}\longrightarrow w_{30}\longrightarrow R_{36}\longrightarrow G_9
\]

ordonne le relèvement ; les \(396\) niveaux, les \(2376\) trits, les \(43\) tours et les \(387\) paires \((K,K+9)\) par canal quantifient le transit ; les cinq régions de \(\pi\) et la fibre ambiante de \(729\) éléments publient des lectures internes ultérieures. Cette séquence rend visible une monodromie productive : le retour retrouve la phase et la coinduction conserve la croissance de mémoire.

\subsubsection{Unité radionale et coordination des lecteurs circulaire, elliptique, projectif et hélicoïdal}

Le radion concentre dans une unité métrologique la structure que le radian publie angulairement. Sa définition naturelle est celle d'une \textbf{section orientée de phase, d'action et de mémoire}. Une unité de phase \(\theta=1\) transporte une unité d'action \(\sigma\hbar_{\mathrm{ret}}\), appartient à un tour complet \(T_+=2\pi_{\mathrm{HMT}}\) et conserve feuillet, orientation, retour nonadique et position dans le cylindre de raffinement.

Le radian exprime la coordonnée angulaire visible. Le radion conserve, avec cette coordonnée, l'action transportée, l'orientation, la retenue et la mémoire du tour. La métrologie angulaire acquiert ainsi une généalogie. Mesurer un arc signifie identifier quelle rotation a été publiée et quel état complet a réalisé cette rotation.

La relation entre phase et aire prend une forme particulièrement claire :

\[
\mathcal A(\theta)=\hbar_{\mathrm{ret}}\,\theta.
\]

Une unité radionale balaie une aire d'action \(\hbar_{\mathrm{ret}}\) ; un tour complet enferme \(h_{\mathrm{ret}}\). Le passage de la phase à l'aire est intégré à l'unité de lecture elle-même. La métrique angulaire, l'action et la mémoire cessent d'apparaître comme des quantités juxtaposées et deviennent les coordonnées d'une même section orientée.

Le radion admet plusieurs publications géométriques coordonnées. Le lecteur circulaire publie la phase ; le lecteur elliptique publie l'action et la forme ; le lecteur intérieur de Klein publie la profondeur ; le lecteur hélicoïdal nonadique publie le retour et la mémoire. L'unité reste généalogiquement identique à travers ces réalisations. C'est pourquoi le radion constitue le lien métrologique approprié entre HMT et la mécanique dimensionnelle : il offre une unité capable de traverser des changements de lecteur tout en conservant l'état que chaque mesure partielle présuppose.

\subsection{Polarité gravitationnelle à cinq composantes et transduction conjointe de l'action, de l'aire, du volume, du temps et de la masse}

La mécanique dimensionnelle permet de définir la dimension comme un régime stable de publication caractérisé par le rang minimal de mémoire indépendante requis par le relèvement. Traverser une dimension signifie appliquer un transducteur qui conserve la généalogie et change le type de lecture. La gravité occupe le niveau de connexion entre ces régimes.

Dans cette formulation, la \textbf{gravitation interdimensionnelle à cinq polarisations corrélatives}, en abrégé \textbf{structure pentapolaire de la gravitation}, est la connexion qui transporte l'unité généalogique et laisse à chaque interface une torsion minimale. Une écriture compacte de cette trace est

\[
\Theta_d^{\min}
=
\nabla_{d+1}\circ T_{d\to d+1}
-
T_{d\to d+1}\circ\nabla_d.
\]

L'expression mesure la différence entre transporter après avoir dérivé et dériver après avoir transporté. Cette différence est la mémoire géométrique minimale du franchissement dimensionnel. La composition d'interfaces produit une holonomie gravitationnelle

\[
\mathfrak G_{d\to D}
=
\operatorname{Hol}
\bigl(
\Theta_d^{\min},
\Theta_{d+1}^{\min},
\ldots,
\Theta_{D-1}^{\min}
\bigr).
\]

L'objet fondamental est cette connexion interdimensionnelle. Un mode local de type graviton occupe le rang dérivé de publication quantifiée de la torsion ; la gravitation conserve sa résidence primaire dans la compatibilité entre dimensions.

Le terme pentapolaire rassemble les cinq polarisations consubstantielles de l'état continu :

\[
\mathbf G_d^{\pm}
=
\bigl(
G_{\mathrm{sp},d}^{\pm},
G_{\mathrm{sol},d}^{\pm},
G_{\mathrm{gau},d}^{\pm},
G_{\mathrm{coh},d}^{\pm},
G_{\mathrm{det},d}^{\pm}
\bigr).
\]

Les cinq composantes émergent corrélativement du même état enrichi. La polarité \(\pm\) enregistre concentration directe et publication conjuguée. La gravité traverse les strates dimensionnelles en transportant simultanément ces cinq lectures ; la torsion minimale conserve la trace de chaque franchissement.

La géométrie acquiert ainsi une image unifiée. Le régime elliptique organise les clôtures ; le parabolique organise les seuils ; l'hyperbolique organise la profondeur. La transduction dimensionnelle relie ces régimes aux réponses radiales \(p\), aux holonomies de spin et aux publications de masse. La structure pentapolaire de la gravitation exprime la cohérence globale de ce lien.

La théorie M fournit des écrans physiques ultérieurs pour cette connexion : dimensions supplémentaires, cordes, branes et dualité T réalisent des modes concrets de transduction. HMT apporte la généalogie APP–TRIT–TPK, l'état enrichi, la mémoire et la lecture bidirectionnelle ; l'interface avec la théorie M publie ces structures dans des géométries physiques typées. La dualité T échange des lecteurs dimensionnels compatibles et conserve la classe généalogique de l'état.

\subsubsection{Compatibilité des réalisations classique, relativiste, quantique et HMT--MD}

La mécanique classique publie trajectoire et quantité de mouvement. La relativité publie structure causale, métrique et courbure. La mécanique quantique publie phase, holonomie, vacances et amplitudes. La mécanique dimensionnelle fournit les transducteurs qui relient ces publications, et HMT conserve la généalogie discrète qui les soutient.

Dans la lecture classique, une orbite décrit la projection visible du mouvement. Dans la lecture relativiste, cette orbite est intégrée à la géométrie du lecteur. Dans la lecture quantique, la phase et le spin enregistrent l'holonomie de la section. Dans la lecture HMT–MD, les trois expressions appartiennent à une même biographie relevée : le radion quantifie phase, action et mémoire ; \(p\) classe la réponse radiale ; \(\tau\) classe la courbure ; l'holonomie publie le spin ; le transducteur dimensionnel publie la masse ; la connexion pentapolaire publie la gravitation.

La dissipation classique apparente acquiert une résidence dans le bilan généalogique élargi. L'accélération relativiste s'exprime par connexion et projection. L'indétermination quantique s'organise par lecteurs, vacances et branches d'Ambivalence. Le principe de Mach relie la réponse locale à la distribution globale. Le Landauer nul garantit la réversibilité informationnelle de l'état complet. Tous ces éléments convergent vers une conservation plus profonde : l'identité de l'unité à travers les changements de publication.

\subsubsection{Théorème narratif de conservation généalogique et de prolongation holographique}

La clôture peut désormais se raconter comme une seule séquence. L'unité se décompose sans perdre son identité. La phase orientée acquiert une mesure dans le radion. Le retour nonadique transforme le cycle en hélice par la mémoire. Le caractère \(p\) transforme l'hélice en familles radiales et classe ses régimes. La persistance de ces trajectoires produit des sections ayant spin, charge, saveur, couleur et masse. Les vacances ouvrent des canaux ; la polarité électrique transporte la phase ; la réponse magnétique courbe transversalement ; les seuils publient le rayonnement. L'autoconservation clôt le bilan généalogique d'action et de quantité de mouvement. L'auto-inertie relationnelle relie connexion propre, projection et totalité machienne. L'Ambivalence conserve les directions directe et conjuguée. La transduction aire–volume conduit à Barbero–Immirzi et à la loi des masses. La torsion minimale enchaîne les dimensions et publie la structure pentapolaire de la gravitation.

La croissance vers l'intérieur et l'expansion vers l'extérieur sont des orientations conjuguées du même processus. Le raffinement hyperbolique génère la profondeur ; la transduction dimensionnelle génère l'extension. La retenue \(0/1\) relie les échelles ; la monodromie relie les tours ; la mémoire relie les états ; l'holonomie relie les orientations ; la gravitation relie les dimensions. Chaque lien conserve une unité qui évolue et chaque publication finie contient la règle de sa prolongation.

La thèse fondamentale est ainsi formulée :

\begin{quote}
\textbf{La clôture holographique de l'infini réalise le principe de conservation évolutive de l'unité par transduction dimensionnelle et projection holographique : l'unité maintient son identité généalogique en se déployant comme phase, action, courbure, mémoire, particule, champ et gravitation ; chaque lecture finie conserve la règle de reconstruction et de prolongation de l'état complet.}
\end{quote}

Cette formulation offre une image physique et mathématique conjointe. L'infini acquiert une loi de clôture locale ; l'unité acquiert de la profondeur ; la dimension acquiert de la mémoire ; la conservation acquiert un caractère relationnel ; la force acquiert l'autoconservation ; l'inertie acquiert une structure machienne ; la gravité acquiert une connexion pentapolaire ; et l'holographie acquiert le sens précis de reconstruction entre lecteurs compatibles.
